% !TEX program = xelatex
\documentclass[a4paper,11pt]{article}
\usepackage[margin=2cm]{geometry}
\usepackage{amsmath,amssymb,graphicx,caption,microtype,indentfirst}
\usepackage[hidelinks]{hyperref}
\graphicspath{{../../random_tfim/results/figures_imaginary_time_fixed/}}
\captionsetup{font=small,labelfont=bf}
\setlength{\parskip}{0.35em}
\title{Imaginary-time Autocorrelation for RTFIM}
\author{Zihao Xu}
\date{October 5, 2026}
\begin{document}
\maketitle
\vspace{-2em}

\section{Random transverse-field Ising model}
We consider a periodic chain of length $L$ with Hamiltonian
\begin{equation}
 H=-\sum_{i=1}^{L}J_i\sigma_i^z\sigma_{i+1}^z
   -h\sum_{i=1}^{L}\sigma_i^x,
 \qquad J_i\sim U(0,1),\quad h=h_0/e.
\end{equation}
The bonds are independent and the transverse field is uniform. The
thermodynamic critical point satisfies $\overline{\ln J}=\ln h$.
Since $\overline{\ln J}=\int_0^1\ln J\,dJ=-1$, it occurs at $h_{0,c}=1$.
For $1<h_0<e$, the system is globally paramagnetic but can contain rare,
locally ordered segments with very small tunnelling gaps. For $h_0>e$,
the field exceeds every bond and the system is a gapped disordered
paramagnet. Thus
\begin{equation}
 \begin{cases}
  \text{Ferromagnetic phase}, & 0<h_0<1,\\
  \text{Paramagnetic Griffiths region}, & 1<h_0<e,\\
  \text{Gapped paramagnetic region}, & h_0>e.
 \end{cases}
\end{equation}
Here $h_0=e$ is the endpoint of the Griffiths region set by the bounded
bond distribution, rather than a second symmetry-breaking transition.

\section{Imaginary-time autocorrelation and numerical results}
We study the zero-temperature autocorrelation at $j=L/2$,
\begin{equation}
 C_j(\tau)=\langle0|e^{\tau H}\sigma_j^z e^{-\tau H}\sigma_j^z|0\rangle,
 \qquad C_{\rm av}(\tau)=\overline{C_j(\tau)}.
 \label{eq:c}
\end{equation}
The results use $50\,000$ disorder realizations for each $(L,h_0)$.
We compare $C=A\tau^{-\alpha}$ and $C=Ae^{-\lambda\tau}$ using equal-weight
least squares in $\ln C$. Both models use the same finite positive data
at $\tau>0$; the displayed floor $10^{-8}$ does not truncate the fits.
The labels $1/z$ and $1/\xi_\tau$ in the figures denote $\alpha$ and
$\lambda$, respectively.

\begin{figure}[htbp]
 \centering
 \includegraphics[width=\linewidth]{figure1_imaginary_time_autocorrelation.png}
 \caption{Average imaginary-time autocorrelation at $L=128$. Colours distinguish fields; solid curves show disorder means, shading indicates one SEM, and dashed curves show the preferred power-law or exponential fit. Legend rates are $1/z=\alpha$ and $1/\xi_\tau=\lambda$. All finite positive data enter the fits; $10^{-8}$ is only the display floor.}
 \label{fig:c}
\end{figure}

\subsection{Field dependence and crossover}
Figure~\ref{fig:c} shows nearly plateau-like behaviour in the ferromagnetic
region. Its very small fitted power-law rates describe slow decay over
the available times, rather than paramagnetic Griffiths exponents.
Just above criticality, at $h_0\simeq1.096$, the curve is approximately
straight on logarithmic axes and is well described by a power law with
$\alpha\simeq0.166$. Curvature becomes more pronounced as $h_0$ increases.
Around $h_0\simeq1.738$ and $2.089$, neither model reproduces the whole
curve well, although the preferred model changes from power law to
exponential. At larger fields the exponential is increasingly favoured;
preference alone, however, does not imply a close fit.

This intermediate regime reflects the coexistence of several decay
scales. Each sample's autocorrelation is a weighted sum of exponentials.
A broad distribution of rare-region gaps can produce a power law after
disorder averaging, but short-time contributions and finite-size effects
bend the curve away from that asymptote. As $h_0$ approaches $e$ from
below, the relevant ordered segments become rarer and may be poorly
represented by a finite ensemble. Fitting early times and very small
late-time tails together can therefore make both simple models inadequate.
These are plausible explanations of the observed curvature, rather than
causes separately verified by the figures.

The accessible data thus evolve from power-law-like to exponential-like
behaviour with increasing $h_0$. This does not establish an exponential
thermodynamic asymptote inside $1<h_0<e$: the true Griffiths tail may
require larger systems, more samples or a later-time fitting window.

\begin{figure}[htbp]
 \centering
 \includegraphics[width=\linewidth]{figure2_imaginary_time_fit_parameters.png}
 \caption{Field dependence of the power-law rate (upper left), exponential rate (upper right), and smaller fit score (lower left), with colours distinguishing system sizes. The lower-right panel compares both models at $L=128$. The scores are unweighted log-space residuals divided by $N-2$; their crossing near $h_0=1.77$ marks a change in model preference, not a phase boundary.}
 \label{fig:scan}
\end{figure}

Figure~\ref{fig:scan} shows increasing fitted decay rates with field.
The minimum residual is largest in the intermediate regime and decreases
toward large fields, consistent with Fig.~\ref{fig:c}. The four sizes
show similar trends, but this agreement does not establish convergence
of the rare-event tail.

\subsection{Weighting and interpretation of the fit scores}
SEM weighting is not used because the aim is to characterize decay over
several time decades. In logarithmic coordinates, propagated SEM weights
are approximately $w=(C/\mathrm{SEM})^2$. For $L=128$, points with
$\tau<1$ carry about $99.40\%$ of the total weight at $h_0\simeq1.738$
and $99.56\%$ at $h_0\simeq2.089$. Such weighting gives little influence
to the long-time behaviour. This dominance depends on field and refers
to log-coordinate weights, not $1/\mathrm{SEM}^2$ in the original
coordinate. Equal log-space weights balance the decay more evenly,
although they do not account for unequal uncertainties or correlations
between times.

The plotted ``reduced $\chi^2$'' is actually the unweighted log-space
squared residual divided by $N-2$. Its smaller value identifies the
better of two approximations on the chosen window; it is not a standard
error-normalized goodness-of-fit statistic. In particular, a value near
one has no special statistical meaning.

At $L=128$, the scores cross near $h_0=1.77$, well below the Griffiths
endpoint $e$. Their relative magnitude therefore cannot determine the
phase boundary: an exponential can win within the Griffiths region,
and a power law can win for a ferromagnetic near-plateau curve.
For this model, $1<h_0<e$ and $h_0>e$ are assigned from the disorder
distribution. The fits indicate how clearly the available data display
the expected decay. Identifying asymptotic behaviour additionally requires
stable long-time slopes and convergence with fitting window, size and
sample number.

\appendix
\section{JW, Majorana covariance and Pfaffian method}
\subsection{Quadratic Hamiltonian and ground-state covariance}

The Jordan--Wigner transformation defines Hermitian Majoranas
\begin{equation}
 a_i=\left(\prod_{k<i}\sigma_k^x\right)\sigma_i^z,
 \qquad b_i=\left(\prod_{k<i}\sigma_k^x\right)\sigma_i^y.
\end{equation}
They obey $\{\gamma_m,\gamma_n\}=2\delta_{mn}$ and give
$\sigma_i^x=i a_i b_i$ and
$\sigma_j^z=(\prod_{k<j}i a_k b_k)a_j$. The Hamiltonian becomes
\begin{equation}
 H_q=-i a^TK_qb,\qquad
 (K_q)_{ii}=h,\quad (K_q)_{i+1,i}=-J_i,\quad
 (K_q)_{1L}=-qJ_L.
\end{equation}
Here $q=-P$, with spin parity $P=\prod_i\sigma_i^x$; bond contributions
are added when matrix entries coincide. For the periodic spin chain, the even ground
state uses antiperiodic fermions ($q=-1$). Inserting $\sigma_j^z$ flips
parity, so the intermediate evolution uses periodic fermions ($q=+1$).
This sector change is essential.

A singular-value decomposition $K_q=U_qS_qV_q^T$, with singular values
$s_{q,\mu}$, determines modes with excitation energies $2s_{q,\mu}$.
The initial vacuum energy is $E_0=-\sum_\mu s_{-,\mu}$.
The Gaussian ground
state is specified by its Majorana covariance matrix,
\begin{equation}
 \Gamma_{mn}=\frac{i}{2}\langle[\gamma_m,\gamma_n]\rangle,
 \qquad \langle\gamma_m\gamma_n\rangle=\delta_{mn}-i\Gamma_{mn}.
\end{equation}
In the ordering $\gamma=(a_1,\ldots,a_L,b_1,\ldots,b_L)$, it is
\begin{equation}
 \Gamma_0=\begin{pmatrix}0&G_0\\-G_0^T&0\end{pmatrix},
 \qquad G_0=U_-V_-^T.
\end{equation}

\subsection{Spin insertion and imaginary-time evolution}

Set $|\psi_j\rangle=\sigma_j^z|0\rangle$. Conjugation by the
Jordan--Wigner string changes signs of Majoranas and preserves Gaussianity.
The inserted state's covariance has off-diagonal block
\begin{equation}
 G_j=D_aG_0D_b,\qquad
 (D_a)_{ii}=\begin{cases}-1&i\leq j\\+1&i>j\end{cases},\quad
 (D_b)_{ii}=\begin{cases}-1&i<j\\+1&i\geq j\end{cases}.
\end{equation}
The autocorrelation can therefore be evaluated as the Gaussian overlap
\begin{equation}
 C_j(\tau)=e^{E_0\tau}
 \langle\psi_j|e^{-\tau H_{+}}|\psi_j\rangle.
\end{equation}

In the periodic-fermion evolution basis, define
\begin{equation}
 B=U_+^TG_jV_+,\qquad
 D=\operatorname{diag}(e^{-2s_{+,\mu}\tau}),\qquad
 \delta E=\sum_\mu(s_{+,\mu}-s_{-,\mu}).
\end{equation}
The code's fast evaluation is the equivalent determinant
\begin{equation}
 C_j(\tau)=e^{\delta E\tau}\det R,\qquad
 R=\frac{I+D}{2}+\frac{I-D}{2}B.
\end{equation}
The energy shift accounts for the different vacuum energies of the two
sectors. When $R$ becomes ill conditioned, the code switches to the
signed Pfaffian evaluation below to reduce cancellation in small tails.

\subsection{Gaussian-overlap Pfaffian}

Wick's theorem expresses an even product of Majoranas as the Pfaffian
of an antisymmetric contraction matrix. On the ring, the implementation
uses its Gaussian-overlap form to retain the change of fermion sector.
Choose a Fock reference with occupations $r_\mu\in\{0,1\}$ that has
nonzero overlap with $|\psi_j\rangle$. Relative particle--hole flips
allow the odd-parity state to be represented by an even Gaussian chart,
\begin{equation}
 S_r=\operatorname{diag}(1-2r_\mu),\qquad
 Z=(I-S_rB)(I+S_rB)^{-1},\qquad Z^T=-Z.
\end{equation}
The code selects the reference by sequentially choosing the more probable
occupation, avoiding a singular choice. Write $d_\mu=e^{-2s_{+,\mu}\tau}$
and define diagonal matrices
\begin{equation}
 (K_r)_{\mu\mu}=\begin{cases}d_\mu&r_\mu=1\\1&r_\mu=0\end{cases},
 \qquad
 (Q_r)_{\mu\mu}=\begin{cases}1&r_\mu=1\\d_\mu&r_\mu=0\end{cases}.
\end{equation}
Then the required $2L\times2L$ antisymmetric matrix and overlap are
\begin{equation}
 W=\begin{pmatrix}Z&-K_r\\K_r&-Q_rZQ_r\end{pmatrix},\qquad
 C_j(\tau)=(-1)^{L(L+1)/2}
 \frac{\operatorname{pf}W}{\det(I+Z)}e^{\delta E\tau}.
\end{equation}
The denominator normalizes the Gaussian state, and the fixed sign follows
from the displayed block ordering. Computing only $\sqrt{\det W}$ would
lose the Pfaffian sign. The implementation tracks the sign and
$\ln|\operatorname{pf}W|$ separately and combines the normalization and
energy shift before exponentiation.

Finally, the mean and SEM are computed over independent disorder
realizations. Checks include $C_j(0)=1$, the decoupled-spin result
$C_j(\tau)=e^{-2h\tau}$, and agreement with direct spin diagonalization
for small periodic chains.
\end{document}

